% se puede agregar la opción [english] para 
%  memorias o tesis en inglés (borrando el archivo .aux)
\documentclass{umemoria} 

\depto{Departamento de ???}
\author{Nombre Completo Autor}
\title{Título de la Memoria/Tesis}

% incluir ambos comandos para una doble titulación
%  o quitar el comando que no aplica
\memoria{Ingenier[o/a?] Civil en ???}
\tesis{Magíster en ???}
%\tesis{Doctor en ???} % incluir solo este comando para doctorados

% puede haber varios profesores guía seperados por coma;
% pero si es una memoria, solo puede haber un profesor guía
\guia{Nombre Completo Guía} 

% puede haber varios profesores co-guía seperados por coma;
% pero si es una memoria, el profesor co-guía será el primer
% integrante de la comisión
%\coguia{Nombre Completo Co-Guía} % incluir en caso de co-guía de *tesis*

%\cotutela{Nombre Institución} % incluir en caso de cotutela

\comision{Nombre Completo Uno,Nombre Completo Dos,Nombre Completo Tres}

%\auspicio{Nombre Institución} % incluir en caso de recibir financiamiento

% tiene que ser el año en que se da el examen de título/grado (defensa)
%\anho{2021} % incluir solo para reemplazar el año actual

\usepackage{lipsum}

\begin{document}

\frontmatter
\maketitle

\begin{resumen}
\lipsum[1-4]
\end{resumen}

% opcional: incluir para tesis en inglés;
%  en este caso hay que tener el resumen y abstract
%   en ambos idiomas
%\begin{abstract}
%\lipsum[1-4]
%\end{abstract}

\begin{dedicatoria}
Una dedicatoria corta.
\end{dedicatoria}

\begin{thanks}
\lipsum[1-2]
\end{thanks}

\tableofcontents
\listoftables % opcional
\listoffigures % opcional

\mainmatter

\chapter{Introducción}

\lipsum[30-35]
\begin{enumerate}
	\item Item 1
	\begin{enumerate}
		\item Subitem 1
		\item Subitem 2 (ver Figura \ref{logofcfm})
	\end{enumerate}
	\item Item 2
	\item Item 3
\end{enumerate}
\begin{teo}
Se tiene que $$\int_0^t e^sds=e^t-1.$$
\end{teo}
\lipsum[36-40]
\chapter{Segundo}

\lipsum[1-3]

\begin{defn}[ver \cite{KAR00}] Definición definitiva $$\frac{d}{dx}\int_a^xf(y)dy=f(x).$$\end{defn}
\chapter{Tercero}
\lipsum[50-60]
\chapter{Conclusión}

\lipsum[130-132]
\begin{figure}
	\centering
	\includegraphics[scale=.2]{imagenes/fcfm.pdf}
	\caption{Logo de la Facultad}
	\label{logofcfm}
\end{figure}

\lipsum[133-134]

\begin{table}
	\centering
	\caption{Tabla 1}
	\label{tabla:1}
	\begin{tabular}{|c|c|r|}
		\hline
		\textbf{Campo 1} & \textbf{Campo 2} & \textbf{Num} \\\hline
		Valor 1a & Valor 2a & 3\\
		Valor 1b & Valor 2b & 3\\
		\hline
	\end{tabular}

\end{table}

\lipsum[135]


% ver https://www.overleaf.com/learn/latex/Glossaries
% 
\newglossaryentry{latex}
{
    name=latex,
    description={Is a mark up language specially suited for scientific documents}
}

\newacronym{cadcc}{CADCC}{Centro de Alumnos del Departamento de Ciencias de la Computación}

\newacronym{dcc}{DCC}{Departamento de Ciencias de la Computación, de la Universidad de Chile}

\newacronym{ps}{CC5402}{CC5402 Proyecto de software} 

\newacronym{e}{CC6908}{CC6908 Introducción al trabajo de título}

\newacronym{f1}{CC6909}{CC6909 Trabajo de Titulo}

\newacronym{f2}{CC6910}{CC6910 Trabajo de Memoria de Titulo}

\newacronym{S}{\textit{Alumno S}}{alumno auto suficiente}

\newacronym{R}{\textit{Alumno R}}{alumno regular}

\newacronym{P}{\textit{Alumno P}}{alumno práctico}

\newacronym{UCD}{UCD}{User Centered Desing}

\newacronym{f}{FG}{\textit{Focus Group}}

\newglossaryentry{i*}{
    name=i*,
    description={El lenguaje de modelamiento i*, fue introducido como un marco de referencia orientado a modelar actores y objetivos. Consiste en un lenguaje de modelamiento junto con una serie de tecnicas que permiten analizar estos modelos. \cite{Dalpiaz2016}}
    }
    
\newglossaryentry{effectiveGUI}{
    name=Efective GUI,
    description= {Cuando se habla de \textit{Efective GUI} o interfaz gráfica efectiva para el usuario, se habla de una interfaz que está diseñada para simular o reemplazar una interacción humana, se acuña principalmente en el área de IA relacionada con los chatbots}
    }

\newglossaryentry{Telegram}{
    name=Telegram,
    description= {Telegram es un software gratiuto, multiplataforma, basado en la nube, de mensajería instantanea. El servicio además provee video encriptado de principio a fin, VoIP, envío de archivos y varias otras funciones. Fue lanzado para iOS en Agosto 14 2013 y Android en Octubre 2013 }
}

\newglossaryentry{Celery}{
    name=Celery,
    description= {Telegram es un software gratiuto, multiplataforma, basado en la nube, de mensajería instantanea. El servicio además provee video encriptado de principio a fin, VoIP, envío de archivos y varias otras funciones. Fue lanzado para iOS en Agosto 14 2013 y Android en Octubre 2013 }
}

\newglossaryentry{Django}{
    name=Django,
    description= {Telegram es un software gratiuto, multiplataforma, basado en la nube, de mensajería instantanea. El servicio además provee video encriptado de principio a fin, VoIP, envío de archivos y varias otras funciones. Fue lanzado para iOS en Agosto 14 2013 y Android en Octubre 2013 }
}

\newglossaryentry{ORM}{
    name=ORM,
    description= {Telegram es un software gratiuto, multiplataforma, basado en la nube, de mensajería instantanea. El servicio además provee video encriptado de principio a fin, VoIP, envío de archivos y varias otras funciones. Fue lanzado para iOS en Agosto 14 2013 y Android en Octubre 2013 }
}

\newglossaryentry{ODM}{
    name=ODM,
    description= {Telegram es un software gratiuto, multiplataforma, basado en la nube, de mensajería instantanea. El servicio además provee video encriptado de principio a fin, VoIP, envío de archivos y varias otras funciones. Fue lanzado para iOS en Agosto 14 2013 y Android en Octubre 2013 }
}

\newglossaryentry{REST}{
    name=REST,
    description={Representational State Transfer, architectural style for distributed hypermedia systems}
}
 % opcional

\nocite{*}
\bibliographystyle{plain}
\bibliography{bibliografia}

% opcional ...
\begin{appendices}
\chapter{Anexo}
\lipsum[50-60]
\end{appendices}
\end{document}
